\documentclass[10pt,a4paper]{amsart}
\usepackage[margin=19mm]{geometry}
\usepackage{amsmath,amssymb,mathtools}
\usepackage[scr=rsfs,cal=euler]{mathalfa}
\usepackage{xfrac}
\usepackage[unicode,hidelinks,bookmarks=false]{hyperref}
\newcommand{\ie}{i.e.\ }
\newcommand{\cf}{cf.\ }
\newcommand{\resp}{respectively\ }
\newcommand{\Subsection}{\subsection}
\providecommand{\nobreakdash}{\nobreak\hbox{-}}
\setlength{\emergencystretch}{2em}
\multlinegap0pt
\usepackage{bm}
\usepackage{bbm}
\usepackage{dsfont}
\usepackage{xspace}
\usepackage{cancel}
\usepackage{mathtools}
\usepackage{cleveref}
\usepackage[shortlabels]{enumitem}
\usepackage{amsthm}
\makeatletter
\@ifundefined{theorem}{\newtheorem{theorem}{Theorem}}{}
\@ifundefined{lemma}{\newtheorem{lemma}[theorem]{Lemma}}{}
\makeatother
\makeatletter
\expandafter\ifx\csname named\endcsname\relax
\newenvironment{named}[1]{\renewcommand{\theoremname}{#1}\begin{namedtheorem}}{\end{namedtheorem}}
\fi
\makeatother
\title[Rigidity of highly twisted plat diagrams]{Rigidity of highly twisted plat diagrams}
\author{Nir Lazarovich, Yoav Moriah, Tali Pinsky, Jessica S. Purcell}
\date{Source: arXiv:2508.07303v1 (CC BY 4.0)}
\begin{document}
\maketitle

\noindent\emph{Source and licence.} Rigidity of highly twisted plat diagrams. Version arXiv:2508.07303v1 (CC BY 4.0), distributed under the Creative Commons Attribution 4.0 International licence, as recorded in the arXiv licence metadata for this exact version and shown on the abstract page https://arxiv.org/abs/2508.07303v1. Full rights notice: licenses/arxiv-2508.07303-CC-BY-4.0.txt.\par\smallskip

\noindent\textit{Editorial scope.} Counted unit: the Lemma whose source label is \texttt{lem:arc one interface twregion neg chi} in the article (source0.tex lines 1262--1265, source label \texttt{lem:arc one interface twregion neg chi}), delivered together with the article's own complete original proof of it (source0.tex lines 1267--1350, one \texttt{proof} environment of 84 lines). No mathematics is written, repaired or re-derived: statement and proof are the article's own text line for line. Required context retained uncounted: menasco assumptions (lines 762--785); Lem:ChiPlusFormula (lines 1008--1019); euler of layer is sum over arcs (lines 1032--1039); finding I arcs (lines 1048--1058); one twist region I-arcs (lines 1049--1057); two twist regions I-arcs (lines 1049--1057). Every definition, display and label of those lines is retained unchanged. Omitted from this package and not redistributed: 1--761, 786--1007, 1020--1031, 1040--1047, 1058--1261, 1266--1266, 1351--2375. Nothing inside a retained excerpt is omitted or altered: every retained body is delivered exactly as the article's lines are written, so no omission receipt of any kind accompanies this package. Citations: the retained excerpts carry no citation command at all, so no bibliography entry of the article is transcribed and no reference is redistributed. Reference adaptation: none was needed and none was made; every reference inside the delivered unit resolves inside it, so no unresolved reference or citation is produced. Printed slips kept literal: no printed word, symbol, hypothesis or number of the counted statement or of its proof was altered. Layout and class: the article's own class is \texttt{amsart} with packages including amsmath, amsthm, amssymb, graphicx, color, import, xy, subcaption, hyperref, cleveref, overpic, enumerate, todonotes, ulem; the wrapper is the lane's pinned \texttt{amsart} editorial wrapper at 10pt on A4, which restates the theorem-like environments the retained text needs and adds the article's own macro definitions that the retained text needs (none), with extra wrapper packages (bm, bbm, dsfont, xspace, cancel, mathtools, cleveref, enumitem). The delivered numbering is the unit's own. Assets: no retained span contains a figure environment, a graphics-inclusion command, a PDF-page inclusion, a table or a TikZ picture; no image file is copied into this package, the package declares no asset dependency and no image file is redistributed with it. Licence: version arXiv:2508.07303v1 of this article is distributed under the Creative Commons Attribution 4.0 International licence, as recorded in the arXiv licence metadata for this exact version and shown on the abstract page https://arxiv.org/abs/2508.07303v1. A full rights notice is delivered at licenses/arxiv-2508.07303-CC-BY-4.0.txt. Characters: every retained excerpt is pure ASCII, so the delivered engine, XeLaTeX with its default Latin Modern text font, reports no missing character.
\begin{enumerate}[label = ($\diamondsuit$\arabic*)]
\item\label{no vertical intersection}  $S\cap P^\pm$ does not intersect the 
vertical segments of $K$.

\item \label{menasco no scc} Every component of $S\cap P^\pm$ meets a twist region or 
an interface point. That is, there are no components that bound disks disjoint from 
$K$ in the diagram. 

\item\label{menasco assumptions} No component of $S\cap P^\pm$ visits a twist 
box more than once.

\item\label{no scc in layers} If a component of $S\cap P^\pm$ is contained in the 
layer $\ell_i$, then it is an O-curve (and $i$ is even).

\item\label{no trivial U-arcs} 
If a component of $S\cap P^\pm\cap \ell_i$, $1\le i\le n$, is a U-arc with at 
least one free endpoint, then it must pass through a 
twist region.

\item\label{no double intersection outside} 
No component of $S\cap P^\pm$ meets a layer in a pair of arcs that both have 
free endpoints. That is, for any curve $c$ of $S\cap P^\pm$, and among all 
arcs of $c\cap \ell_i$, at most one such arc has any free endpoint(s).
\end{enumerate}

\begin{lemma}\label{Lem:ChiPlusFormula}
Equivalently,
\begin{equation}
    \chi_+(\alpha)=\begin{Bmatrix}
    1 & \text{if }\alpha\text{ is an O-curve}\\
    0 & \text{if }\alpha\text{ is an extremal U-arc}\\
    -\tfrac12 & \text{otherwise}
    \end{Bmatrix} + \tfrac14\#\{\text{free endpoints of }\alpha\}
\end{equation}
The only components $\alpha$ of $S\cap P^{\pm}\cap \ell_i$ with positive 
$\chi_+(\alpha)$ are O-curves or extremal U-arcs with at least one free endpoint.
\end{lemma}

\begin{lemma}\label{euler of layer is sum over arcs}
The value of $\chi_i$ satisfies
    \begin{equation}\label{eq: euler by arc}
        \chi_i = \sum_{\alpha} \chi_+(\alpha) 
        \; -\;\#\{\text{internal saddles in }\ell_i\}
    \end{equation}
    where the sum runs over every component $\alpha$ of  $S\cap P^\pm \cap \ell_i$.
\end{lemma}

\begin{lemma}\label{finding I arcs}\hfill
\begin{enumerate}[label=(\arabic*)]
\item\label{one twist region I-arcs} If $S \cap P^\pm \cap \ell_i$ meets the 
twist region $T$ but does not meet its neighboring twist regions, then there 
are at least two I-arcs through $T$. 
\item\label{two twist regions I-arcs} If $S \cap P^\pm \cap \ell_i$ meets two 
neighboring twist regions $T,T'$ but does not meet their neighboring twist 
regions, and there is at most one arc in $S\cap P^+$ connecting $T$ and $T'$, 
then there are at least two I-arcs through $T$ or $T'$.
\end{enumerate}
\end{lemma}

\begin{lemma}\label{lem:arc one interface twregion neg chi}
Assume there are no O-curves and there exists an arc $\alpha$ with at most one 
interface endpoint that passes through a twist region. Then $\chi_i<-1.$
\end{lemma}

\begin{proof}
The arc $\alpha$ must have at least one free endpoint, so $i$ is even. 
Enumerate the twist boxes $T_1,\dots,T_m$, and denote by $k_1,\dots,k_m$, $s_1,\dots,s_m$ 
the number of crossings in $T_1,\dots,T_m$ and the number of components of 
$S\cap T_1,\dots,S\cap T_m$ respectively.

Since $\alpha$ has at least one free endpoint, it must (a) pass through $T_1$ and have a 
free endpoint to its left, or (b) pass through $T_m$ and have a free endpoint to its 
right (or both). Let $F_L$ be the number of free endpoints corresponding to (a), and 
$F_R$ be the number of free endpoints corresponding to (b). Note that each free endpoint 
is actually counted twice, once as an endpoint of a curve in $P^+$ and once in $P^-$. 
Then $F_L \le 2k_1 \cdot s_1$ and $F_R \le 2k_m\cdot s_m$. Let us denote by 
$G_L = 2k_1s_1 - F_L$ (resp.\ $G_R = 2k_m s_m - F_R$) the arcs meeting the left 
side of $T_1$ (resp.\ the right side of $T_m$) and not ending in a free endpoint. 
    
Note that 
\begin{multline}
\tfrac14 F_L  - \#\{\text{internal saddles of }T_1\} = \tfrac14 (2k_1 \cdot s_1) - 
\tfrac14 G_L  - (k_1-2)s_1 
\\
= (-\tfrac12 k_1 +2)s_1 - \tfrac14 G_L \le - \tfrac14 G_L     
\end{multline}
where the last inequality follows from the assumption $k_1\ge 4$.
Similarly, the assumption $k_m\ge 4$ implies
\[
\tfrac14 F_R  - \#\{\text{internal saddles of }T_m\} \le - \tfrac14 G_R
\]

We need to compute $\chi_i = \sum \chi_+(\alpha) - \#\{\text{internal saddles in }\ell_i\}$. 
By \Cref{euler of layer is sum over arcs}, \Cref{Lem:ChiPlusFormula} and the fact that 
there are no O-curves, this is given by
\begin{align*}
\chi_i &= \sum \chi_+(\alpha) - \#\{\text{internal saddles in }\ell_i\} \\
&= \tfrac14 \#\{\text{free endpoints of arcs}\} - \tfrac12 \#\{\alpha\text{ non-extremal}\} - \#
\{\text{internal saddles in }\ell_i\} \\
&= \tfrac14 (F_L + F_R) -\tfrac12\#\{\alpha\text{ non-extremal}\} - \#\{\text{internal 
saddles of }T_1,T_m\} - \\
& \hspace{3cm}  \#\{\text{internal saddles of }T_2,\dots, T_{m-1}\}
\end{align*}
The above inequalities relating $F_L, F_R$ and $G_L, G_R$ imply
\begin{equation}\label{Eqn:ChiBound}
\chi_i \leq -\tfrac14 (G_L + G_R) - \tfrac12\#\{\alpha\text{ non-extremal}\} 
- \#\{\text{internal saddles of }T_2,\dots, T_{m-1}\}
\end{equation}

\textbf{Case 1.} Suppose some $s_j \in \{s_2, \dots, s_{m-1}\}$ is nonzero. Then by the 
assumption that $k_j\geq 4$, there are at least two internal saddles. In this case the 
third term in the inequality \eqref{Eqn:ChiBound} is then at most $-2$, with the other 
two terms non-positive, so we have $\chi_i \leq -2 < -1$. 

\textbf{Case 2.} If all of $s_2= \dots = s_{m-1} =0$, \Cref{Eqn:ChiBound} becomes
\begin{equation}\label{Eqn:ChiBound2}
\chi_i \leq -\tfrac14 (G_L+G_R) - \tfrac12\#\{\alpha\text{ non-extremal}\}
\end{equation}
We divide into further cases.

Consider an arc in $S\cap P^+\cap \ell_i$ that runs from the left side of $T_1$, around 
the annulus $\ell_i$ avoiding all other twist regions, to connect to the right side of $T_m$. 
The existence of any such arc implies there is another such arc on $P^-$. These are each 
counted in both $G_L$ and $G_R$. Thus any such arc contributes $4$ to $G_L+G_R$.  
    
\textbf{Case 2.A.} There are at least two such arcs in $S\cap P^+\cap \ell_i$. 
Then the first term of \Cref{Eqn:ChiBound2} is at most $-\tfrac14\cdot 8 \leq -2 < -1$, 
and the second is nonpositive, so we have the result in this case. 

\textbf{Case 2.B.} There is exactly one such arc in $S\cap P^+\cap \ell_i$. Then 
\Cref{finding I arcs}~\ref{two twist regions I-arcs} implies that there are (at least) 
two I-arcs. This means that the first term of \Cref{Eqn:ChiBound2} is at most 
$-\tfrac14\cdot 4$ and the second is at most $-\tfrac12 \cdot 2$, so the sum is at most
\[ \chi_i \leq -1 - 1 < -1 \]

\textbf{Case 2.C.}
Suppose there are no such arcs. The case $s_1=s_m=0$ does not occur because of our 
assumptions, so suppose without loss of generality that $S\cap P^\pm \cap \ell_i$ 
meets $T_1$. Then, by \Cref{finding I arcs} \ref{one twist region I-arcs}, there 
are at least two I-arcs $\alpha_1$, $\alpha_2$ passing through $T_1$.

Each of $\alpha_i$ belongs to a simple closed curve of $S\cap P^\pm$ that intersects 
$\ell_i$ in another I-arc, say $\beta_i$. The curves $\alpha_1$, $\alpha_2$, $\beta_1$,
$\beta_2$ must all be distinct by \Cref{menasco assumptions}. This means we have four 
distinct I-arcs in this case. 
Then the second term of \Cref{Eqn:ChiBound2} is at most $-\tfrac12 \cdot 4 = -2$, and 
the first term is non-positive, so $\chi_i < -1$. 
\end{proof}

\end{document}
